\subsection{Elliptic periods from the actual normal flow germs}

\begin{statement}{C.2.UN.29}{Proved}{Rank-one elliptic component covers have only smooth multiple fibres}
Let $E$ be an actual multiplicity-one covering component
constructed in C.2.UN.16, with its saturated two-dimensional
space of limiting translation fields. In the generic-rank-one
branch, take the actual graph resolution and Stein factor
of their ratio. When its connected general fibre is elliptic,
the resulting relatively minimal elliptic surface
\[
                         q:E_0\longrightarrow B
                                                               \tag{C.RG.1}
\]
has identity integral $H_1$ monodromy. Every fibre is a
positive multiple of a smooth elliptic curve, and all smooth
fibres are torsors under one fixed elliptic curve. Moreover
\[
 R^1q_*\OO_{E_0}\simeq\OO_B,\qquad
 p_g(E)=g(B)\le1,\qquad \chi(\OO_E)=0.
                                                               \tag{C.RG.2}
\]
The maps to $E_0$ are the specified point blowups resolving
the first integral followed by vertical exceptional-curve
contractions. No extension of that first integral to an
ambient neighbourhood is assumed.

Together with C.2.UN.22, this proves that every such root-cover
component of nonnegative Kodaira dimension has coherent
Euler characteristic zero. In particular its minimal model
is neither K3 nor Enriques. Every original component of
Kodaira dimension one has an isotrivial elliptic Iitaka
fibration over a curve of genus at most one; its possible
singular fibres are retained by the actual cyclic quotient
comparison, and are not asserted to be smooth.
\end{statement}
\begin{proof}
We work on the actual root-change and resolution model
$f:U\to\Delta_s$ in C.2.UN.16. Its smooth general fibre
is still a complex two-torus. The chosen component $E$
has multiplicity one. Its two ambient meromorphic fields
$V_1,V_2$ are vertical, commute, are holomorphic and
independent on each general torus, and have no pole along
the generic point of $E$. Their restrictions are independent
over $\C$ and have rank one over the meromorphic function
field of $E$.

Resolve the meromorphic ratio on $E$ by point blowups
and take its connected-fibre Stein factor $q:E\to B$.
Perform the same point blowups in the smooth ambient
threefold. This gives the actual ambient extension of
these modifications, retains multiplicity one on the
strict transform of $E$, and lifts both fields
meromorphically. Their generic restrictions and bracket
are unchanged through the common open isomorphism.
We use the same letter $E$ for this modified component.
There is a nonconstant meromorphic function $H$ on $B$
such that
\[
                    V_2|_E=(q^*H)V_1|_E.
                                                               \tag{C.RG.3}
\]
This is precisely the function and Stein map given by
the original ratio, rather than a chosen ambient extension
of it.

Remove finitely many values of $B$ so that $q$ has smooth
connected elliptic fibres and $V_1|_E$ is a nonzero
holomorphic translation field on each of them. The latter
assertion follows from the compact-curve time-form argument
in C.2.UN.16: a semicomplete meromorphic field on an
elliptic curve has no pole. After excluding the finitely
many fibres on which its restriction is identically zero,
it is nowhere zero on every remaining fibre. Remove also
the poles and critical points of $H$.

Let $D$ be the union on $E$ of its intersections with
the other central divisors and with the actual ambient
polar divisors of the two fields. Each remaining general
elliptic fibre meets $D$ in finitely many points, after
discarding any finitely many fibres contained in $D$.
The fields are holomorphic on $U^*$; hence all their
ambient polar divisors lie over zero. In particular
outside this explicitly removed set the ambient fields
and the submersion $f$ are holomorphic. At a suitable
point $p$ of a remaining fibre $F_{b_0}$ there are
holomorphic coordinates $(u,z,s)$ with
\[
 f=s,\qquad V_1=\partial_u,\qquad
 V_2=h(z,s)\partial_u+A(z,s)\partial_z,
 \qquad h(z,0)=H(z).
                                                               \tag{C.RG.4}
\]
Here $z$ is a local coordinate on $B$ along $E$.
Indeed the holomorphic flow-box theorem gives $V_1=\partial_u$,
verticality preserves $s$, and commutation makes both
coefficients of $V_2$ independent of $u$. The local coordinate
$z$ is only used in this one flow box. It need not extend
around the compact elliptic fibre.

Since the ambient fields are independent on the nearby
tori, $A$ is not identically zero. Since their restrictions
to $E$ have rank one, $A(z,0)=0$. Choosing $p$ away
from the zeros of its leading coefficient gives
\[
 A(z,s)=s^d\bigl(a(z)+O(s)\bigr),\qquad d\ge1,
 \qquad a(z)\ne0,\qquad H'(z)\ne0
                                                               \tag{C.RG.5}
\]
on a smaller disc. The order $d$ is finite; it is the
order of the ambient wedge along this component in the
chosen nonvanishing frame.

\par\medskip\noindent\textbf{Actual elliptic return maps.}
Let $\lambda_1(z),\lambda_2(z)$ be the periods of the
time form of $V_1|_{F_z}$, for an oriented integral
homology basis continued over a small disc about $b_0$.
These are holomorphic functions; a local holomorphic
section and the ordinary integration construction of
an elliptic curve supply the marked periods. Put
$c_i=\lambda_i(b_0)$.

The time covering $\C\to F_{b_0}$ of the translation
field pulls the finite bad set $D\cap F_{b_0}$ back to
a discrete subset of $\C$. Choose a piecewise smooth
time path from $0$ to $c_i$ avoiding this discrete set.
Its image is a compact path on which both ambient fields
are holomorphic and $V_1$ is nonzero. A finite chain of
holomorphic flow charts therefore continues the solution
along this entire path for every initial point in a
sufficiently small ambient neighbourhood of $p$.
Its endpoint is again near $p$, since the central path
closes. This constructs a genuine holomorphic return
germ $\Phi_i$; it does not deform a compact elliptic
orbit or assert that such an orbit exists on a nearby
torus.

Every local continuation of the flow of $V_1$ preserves
$f,V_1,V_2$, by verticality and the bracket identity.
Thus the return germ in the single flow box has the form
\[
 \Phi_i(u,z,s)=\bigl(u+\delta_i(z,s),\phi_i(z,s),s\bigr).
                                                               \tag{C.RG.6}
\]
On $E$ the first integral is constant along its elliptic
orbits, and comparison with the continued period gives
\[
 \phi_i(z,0)=z,\qquad
 \delta_i(z,0)=c_i-\lambda_i(z).
                                                               \tag{C.RG.7}
\]
To see the second equality directly, follow the fixed
time path of total time $c_i$ on $F_z$, then return
to the local section by the small flow segment of
time $-\delta_i(z,0)$. The resulting closed path is
the continued integral cycle, so its time integral is
$c_i-\delta_i(z,0)=\lambda_i(z)$. This also specifies
which integral cycle and period enter the return map.

Pushing $V_2$ forward by (C.RG.6) and equating it to
$V_2$ at the endpoint gives the two exact identities
\[
 \begin{split}
 A(\phi_i(z,s),s)&=A(z,s)\,\partial_z\phi_i(z,s),\\
 h(\phi_i(z,s),s)&=h(z,s)+A(z,s)\,\partial_z\delta_i(z,s).
 \end{split}                                                   \tag{C.RG.8}
\]
These are identities of holomorphic germs on the actual
ambient threefold.

\par\medskip\noindent\textbf{The leading normal coefficients.}
The second identity of (C.RG.8), together with $H'$ being
a unit and (C.RG.5), first gives
$\phi_i-z=O(s^d)$. For instance, factor
$h(\phi_i,s)-h(z,s)=(\phi_i-z)h_1(z,s)$;
the factor $h_1$ is a unit because $h_1(z,0)=H'(z)$.
Write
\[
                \phi_i=z+s^d\psi_i(z)+O(s^{d+1}).
                                                               \tag{C.RG.9}
\]
In the first identity of (C.RG.8), all coefficients
of $A$ below order $2d$ cancel. At order $2d$ the
remaining equality is
\[
                         a'\psi_i=a\psi_i'.
                                                               \tag{C.RG.10}
\]
Terms $s^{d+k}a_k(z)$ with $k\ge1$ first contribute
to the difference at order $2d+k$, so they do not
enter this calculation. Since $a$ is a unit,
$(\psi_i/a)'=0$. Hence $\psi_i=k_i a$ for a constant
$k_i\in\C$. At order $d$, the second identity gives
\[
 H'\psi_i=a\,\partial_z\delta_i(z,0)
          =-a\lambda_i',\qquad
             \lambda_i'= -k_i H'.
                                                               \tag{C.RG.11}
\]
Consequently there are constants $d_i\in\C$ with
\[
                         \lambda_i=d_i-k_iH
                    \quad(i=1,2)
                                                               \tag{C.RG.12}
\]
on the small disc. This uses only the two return germs
and one flow box. In particular it does not require
an ambient extension of $q$ or of $H$ around a fibre.
The formal coefficient extraction is independently
certified in \path{code/A_return_germs.py}; its executed
output records the higher-coefficient cancellations.

\par\medskip\noindent\textbf{Global periods and the singular fibres.}
Let $B^\circ$ be a connected complement of finitely
many values containing the disc used above, on which
the elliptic fibration, time form, and $H$ are regular.
Continue the chosen homology basis and periods on its
universal cover. Both sides of (C.RG.12) are holomorphic
there, so the identity theorem gives the same formula
everywhere, with the same constants. Since $H$ is an
actual single-valued meromorphic function on $B$, each
continued period is single-valued on $B^\circ$.
The integration map from $H_1(F_b,\Z)$ to $\C$ is
injective for the time form of a nonzero elliptic
translation field. Thus returning both periods
unchanged forces every integral monodromy matrix to
be the identity.

The oriented ratio $\tau=\lambda_2/\lambda_1$ is
therefore a single-valued holomorphic map
$B^\circ\to\mathbb H$. The denominator never vanishes,
because a nonzero integral elliptic cycle has nonzero
period. The Cayley transform $(\tau-\mathrm i)/(\tau+\mathrm i)$
is bounded, extends holomorphically across all finitely
many omitted points, and is constant on the compact
curve $B$. Hence $\tau$ is constant. Integrating a
locally normalized time form identifies the smooth
fibres with torsors under one fixed elliptic curve.

Contract the vertical exceptional curves to obtain
(C.RG.1). The integral monodromy is unchanged. The
exact Kodaira-table and multiple-root comparison in
(C.RF.3) applies: for a full fibre $mC$, write
$t=u h^m$, make the actual change $t=s^m$, and
normalize by $w^m=u$, $s=wh$. This covering surface
is \'etale over the original one. Its reduced
relatively minimal fibre has identity monodromy,
so is the smooth type $I_0$. It follows that every
original reduced support $C$ is smooth elliptic.
This is precisely the comparison established from
\cite[Theorem~6.2 and Table~I]{Kodaira1963II} in
C.2.UN.22; it includes all multiple fibres.

\par\medskip\noindent\textbf{The direct image and the genus bound.}
Over a small smooth base disc, a local holomorphic
section and integration of the marked differential
give an actual product with the fixed elliptic curve.
Since the integral monodromy is the identity, the
origin-preserving linear part of every overlap map
is the identity. The overlap maps are translations,
which act trivially on $H^1(\OO)$.
At a multiple value use the actual normalized root
family just described. It is smooth with constant
elliptic modulus. The cyclic action on its central
elliptic fibre is free, since the map to $E_0$ is
\'etale, and is therefore a translation. Its linear
part on the neighbouring marked elliptic fibres
is the same discrete automorphism and hence is
the identity. On $H^1(\OO)$ it acts trivially.
Taking invariants in the local root disc consequently
gives $R^1q_*\OO=\OO$ at the multiple value as well:
the coefficient functions invariant under
$s\mapsto\zeta_m s$ are exactly the holomorphic
functions of $t=s^m$. These local generators agree
with the preceding translation-invariant generators
over the punctured discs. This proves the first
assertion of (C.RG.2) with its actual descent map.

Connected fibres give $q_*\OO_{E_0}=\OO_B$.
The coherent Leray sequence on the curve now gives
\[
 H^2(E_0,\OO_{E_0})\simeq H^1(B,\OO_B),\qquad
 h^1(E_0,\OO_{E_0})=g(B)+1.
                                                               \tag{C.RG.13}
\]
Serre duality and birational invariance give
$p_g(E)=g(B)$ and $\chi(\OO_E)=0$.
The independent compact-core inequality
C.2.UN.28 gives $p_g(E)\le1$, and therefore
$g(B)\le1$.

Every root-cover component of nonnegative Kodaira
dimension has either generic rank two, covered by
C.2.UN.22, or rank one. In the latter case a rational
general first-integral fibre would make the surface
uniruled and give Kodaira dimension $-\infty$ by
the actual adjunction argument in C.2.UN.20.
Its general fibre is therefore elliptic and the
present proof applies. This proves the stated
coherent-Euler and K3/Enriques exclusions on the
covering components.

Finally let an original component have Kodaira
dimension one. Its actual generically finite
component cover has Kodaira dimension at least
one by pluricanonical pullback, and at most one
by C.2.UN.15. Its elliptic fibration is therefore
its connected-fibre Iitaka fibration. The cyclic
stabilizer acts on this fibration, since it acts
on the pluricanonical field, and the proper
quotient and graph comparisons in C.2.UN.19--20
give the descended fibration on a smooth
birational model of the original component.
Its base is the quotient of $B$, so has genus
at most one. On a general fibre, the subgroup
fixing a general base point must act by
translations: an origin-preserving nontrivial
linear part would give a rational general quotient
fibre. Its normal line is trivial, so adjunction
restricts every positive pluricanonical line to
$\OO_{\mathbb P^1}(-2m)$; all pluricanonical sections
would then vanish on the dense family of these
fibres, contrary to Kodaira dimension one. The finite translation
subgroup of the fixed elliptic curve is locally
constant, because it consists of points of
bounded finite order. Thus the general quotient
fibre is one fixed elliptic curve as well.
This proves isotriviality through the actual
finite quotient, while leaving its exceptional
and singular fibres in place.
\end{proof}
